\documentclass[11pt]{article}
\usepackage[margin=25mm]{geometry}
\usepackage{fontspec}
\setmainfont{Latin Modern Roman}
\usepackage{amsmath,amssymb,mathtools}
\usepackage{unicode-math}
\setmathfont{Latin Modern Math}
\usepackage{xurl,hyperref}
\hypersetup{colorlinks=true,urlcolor=blue}
\setcounter{secnumdepth}{0}
\setlength{\parindent}{0pt}
\setlength{\parskip}{5pt}
\setlength{\emergencystretch}{3em}
\newcommand{\sref}[2]{\hyperlink{source:#1:#2}{[S#2]}}
\newcommand{\cataloguescope}[1]{\paragraph{Recorded target.}#1}
\newcommand{\R}{\mathbb R}
\newcommand{\C}{\mathbb C}
\newcommand{\N}{\mathbb N}
\newcommand{\Z}{\mathbb Z}
\newcommand{\Q}{\mathbb Q}
\newcommand{\F}{\mathbb F}
\newcommand{\eps}{\varepsilon}
\begin{document}
\section{Hypergroup Correspondence for Discrete Conformal Subnets}\label{statement}

\subsection{Definitions and mathematical statement}

\paragraph{Shared notation.}$\mathbb N=\{1,2,\ldots\}$; the standard number fields have their usual meanings. All additional conventions are stated locally.


\paragraph{Definitions and conventions.}
A compact hypergroup generalizes a compact group by allowing the product of two points to be a probability measure, with compatible convolution, identity and involution. Hypergroup elements can act through quantum operations on a net. Strong additivity means removing an interior point of an interval does not reduce the algebra generated by its two remaining subintervals. A discrete subnet is an inclusion with the source\textquotesingle s discrete decomposition condition, which must be formulated even when the index is infinite.

\paragraph{Target retained from the canonical source catalogue.}
\paragraph{Catalogue wording retained for traceability.}
For a strongly additive conformal net $\mathcal A$, establish a bijection between discrete subnets $\mathcal B\subseteq\mathcal A$ and compact hypergroups $K\subseteq\operatorname{QuOp}(\mathcal A)$, including intrinsic definitions that cover the infinite-index case.

\paragraph{Source scope and qualification.}
The compactness, discreteness and action definitions are themselves part of the target. A finite-index correspondence alone does not establish the requested bijection and intrinsic infinite-index framework.
\subsection{Short English statement}
Match discrete conformal subnets with compact hypergroups of quantum operations, including a meaningful infinite-index theory.
\subsection{Research attempt}
\paragraph{Model, date and investigated scope.}
OpenAI GPT-6 Astra (configured model \texttt{gpt-6-astra}), ultra reasoning; generated 11 October 2026. Exhibited distinct faithful cyclic matrix actions with the same fixed algebra and expectation; diagnosed missing reconstruction hypotheses without counterclaiming against conformal-net correspondence.

\paragraph{Sources and literature check.}
Read the supplied definition, catalogue wording, source qualifications and reference entries for this question. The historical source/context and inherited status are retained; no fresh literature or novelty certification is claimed. No external theorem is used without proof in this attempt; explicit model assumptions are distinguished from facts about the source model.

\paragraph{A fixed-point reconstruction ambiguity.} On \(M_2(\mathbb C)\), let \(U_m=\operatorname{diag}(1,e^{2\pi i/m})\), \(m\ge2\), and let the cyclic group act by \(\alpha_j(x)=U_m^jxU_m^{-j}\). Averaging gives
\[
 E_m\begin{pmatrix}a&b\\c&d\end{pmatrix}
 =\begin{pmatrix}a&0\\0&d\end{pmatrix},
\]
because \(\sum_{j=0}^{m-1}e^{2\pi ij/m}=0\). The fixed algebra and its averaging expectation are therefore identical for every \(m\), while the acting cyclic automorphism groups have different orders. Each action is faithful, since \(\alpha_j\) fixes the off-diagonal matrix unit only when \(j\) is divisible by \(m\).
\paragraph{What fails in a naive bijection.} Recovering an arbitrary compact group merely from the range of an averaging projection cannot be injective in this diagnostic class. Groups are special hypergroups, so adding the word hypergroup cannot by itself repair that loss of action information. Additional canonicality and admissibility hypotheses must do the work.
\paragraph{Source scope.} These matrix actions are inner and are not strongly additive conformal nets or the source's discrete subfactor setting. Thus they do not disprove the proposed source correspondence. They test an incomplete reconstruction argument. The first gap is an intrinsic discrete-inclusion definition and canonical operation-space reconstruction that both exclude irrelevant actions and persist at infinite index. Surjectivity, injectivity, topology and hypergroup convolution in that setting all remain unproved here.

\paragraph{Verification and adversarial audit.}
Computed the phase average and faithfulness; explicitly excluded interpreting inner matrix actions as admissible strongly additive net examples. This is a GPT-6 Astra Ultra self-audit, not independent human verification. Prompt A drafting followed by prompt B scope/dependency/boundary self-audit. Computed the phase average and faithfulness; explicitly excluded interpreting inner matrix actions as admissible strongly additive net examples. The remaining gap is: Define intrinsic infinite-index discreteness and reconstruct a canonical compact hypergroup with provable bijectivity. Unresolved status and zero original-target completion are retained.

Coordinator checked the cyclic averaging expectation and faithfulness, with inner matrix actions distinguished from the source inclusion class. This additional check was model-assisted, not independent human verification.

\paragraph{Tools, process and contributions.}
Python 3 and local file-reading/writing tools for manuscript serialization and structural checks; mathematical calculations were symbolic reasoning, with no numerical experiment or theorem prover. Contribution prompt A was followed by prompt B. GPT-6 Astra with ultra reasoning drafted the research and performed its stated self-audit; any additional coordinator contribution is identified in the verification and audit. A GPT-6-based Codex coordinator assembled the manuscript and checked repository integration; its exact variant and reasoning configuration were not exposed. No human mathematical contribution, independent human proof check or formal proof certification is supplied. Complete derivations and the consulted source references are included here; private reasoning transcripts are not published.

\paragraph{Final status and first remaining gap.}
\textbf{UNRESOLVED}. Define intrinsic infinite-index discreteness and reconstruct a canonical compact hypergroup with provable bijectivity.

\section{Completion Estimate}
\noindent\textbf{COMPLETION: 0\%}
The proved calculation is an elementary auxiliary result or a test of a restricted strategy; no new portion of the full source target is established.

\subsection{Sources}
\begin{itemize}
\item[\hypertarget{source:statement:1}{S1}] ULAM.AI and the UnsolvedMath Contributors, UnsolvedMath, catalogue entry. Catalogue curation under CC BY 4.0; definitions and English exposition added locally.
\url{https://huggingface.co/datasets/ulamai/UnsolvedMath}.
\item[\hypertarget{source:statement:2}{S2}] Subfactors and Applications (2019), Oberwolfach Reports, DOI 10.4171/owr/2019/49, pp. 3053--3140. Primary report accessed 2026-10-11; relevant PDF pages 29, 30, 35, 36.
\url{https://doi.org/10.4171/owr/2019/49}.
\item[\hypertarget{source:statement:3}{S3}] Bischoff--Del Vecchio--Giorgetti, quantum operations on conformal nets (catalogue-listed context; not independently consulted in this preparation).
\url{https://arxiv.org/abs/2204.14105}.
\item[\hypertarget{source:statement:4}{S4}] Bischoff--Del Vecchio--Giorgetti, compact hypergroups from discrete subfactors (catalogue-listed context; not independently consulted in this preparation).
\url{https://arxiv.org/abs/2007.12384}.
\end{itemize}
\label{end:statement}
\end{document}
